\subsubsection{主极点留数与约化行列式：去 Perron 因子的常数层谱指纹}\label{subsec:finite-part-reduced-determinant-residue}
本小节把 Abel/Hadamard 有限部常数分解中出现的主极点留数项 \(C(\theta)\)（见 \ref{def:abel-finite-part} 与命题 \ref{prop:abel-finite-part-mobius-zeta}）
化为一个纯谱对象：\emph{去 Perron 因子后的约化行列式}。
该表述在常数层上提供一个与 moment 阶数兼容的“乘法偏差”指纹，从而可将常数漂移的差异直接还原为有限谱的块间不一致。

\begin{proposition}[主极点留数的谱乘积公式]\label{prop:finite-part-residue-reduced-determinant}
设 \(A\) 为原始（primitive）非负矩阵，Perron 根为 \(\rho=\rho(A)>0\)，并定义
\(
\zeta_A(z):=\det(I-zA)^{-1}
\)。
记主极点 \(z_\star:=1/\rho\)，并令
\[
C(A):=\lim_{z\to z_\star}(1-\rho z)\,\zeta_A(z)
\]
为主极点留数常数（简单极点情形下存在且 \(C(A)>0\)）。
设 \(A\) 的全部特征值按代数重数为 \(\rho=\mu_1,\mu_2,\dots,\mu_d\)。
则有显式闭式
\[
\boxed{\;
C(A)=\prod_{j=2}^{d}\left(1-\frac{\mu_j}{\rho}\right)^{-1}\; }.
\]
等价地，若定义\textbf{约化行列式}
\[
\det\nolimits'\!\left(I-\frac{A}{\rho}\right)
\;:=\;
\prod_{j=2}^{d}\left(1-\frac{\mu_j}{\rho}\right),
\]
则
\[
\boxed{\;
C(A)=\det\nolimits'\!\left(I-\frac{A}{\rho}\right)^{-1}\; }.
\]
\end{proposition}
\begin{proof}
由特征值分解
\[
\det(I-zA)=\prod_{j=1}^{d}(1-z\mu_j)=(1-\rho z)\prod_{j=2}^{d}(1-z\mu_j),
\]
从而
\[
\zeta_A(z)=\frac{1}{1-\rho z}\cdot \prod_{j=2}^{d}(1-z\mu_j)^{-1}.
\]
令 \(z\to z_\star=1/\rho\) 得
\(
C(A)=\prod_{j=2}^{d}(1-\mu_j/\rho)^{-1}
\)。
约化行列式的表述仅为记号改写。证毕。
\end{proof}

\begin{corollary}[moment 阶的约化行列式比值指纹]\label{cor:finite-part-moment-anomaly-reduced-determinant-ratio}
对任意整数 \(q\ge 2\)，令 \(A^{(q)}_\theta\) 表示某个 \(q\)-重 moment-kernel 的原始 realization（参数 \(\theta\) 固定），并记
\(
\rho_q=\rho(A^{(q)}_\theta)
\)、
\(
z_{\star,q}=\rho_q^{-1}
\)，以及其主极点留数常数
\[
C_q(\theta):=\lim_{z\to z_{\star,q}}(1-\rho_q z)\,\det(I-zA^{(q)}_\theta)^{-1}.
\]
定义常数层的 moment 偏差（去指数主尺度后的乘法偏差）为
\[
\mathsf{Anom}^{(q)}_{\mathrm{res}}(\theta):=\log C_q(\theta)-q\log C_1(\theta).
\]
则
\[
\boxed{\;
\mathsf{Anom}^{(q)}_{\mathrm{res}}(\theta)
=
-\log\det\nolimits'\!\left(I-\frac{A^{(q)}_\theta}{\rho_q}\right)
\;+\;
q\,\log\det\nolimits'\!\left(I-\frac{A^{(1)}_\theta}{\rho_1}\right)\; }.
\]
从而其指数化满足约化行列式比值闭式
\[
\boxed{\;
\exp\bigl(\mathsf{Anom}^{(q)}_{\mathrm{res}}(\theta)\bigr)
\;=\;
\frac{\det\nolimits'\!\left(I-\frac{A^{(1)}_\theta}{\rho_1}\right)^{q}}
{\det\nolimits'\!\left(I-\frac{A^{(q)}_\theta}{\rho_q}\right)}\; }.
\]
\end{corollary}
\begin{proof}
对 \(A^{(q)}_\theta\) 与 \(A^{(1)}_\theta\) 分别应用命题 \ref{prop:finite-part-residue-reduced-determinant}，
并取对数相减即得。证毕。
\end{proof}

\begin{proposition}[核版 RH 给出的留数常数压扁界：维数 \texorpdfstring{$\times$}{x} 平方根缺口]\label{prop:finite-part-residue-constant-rh-squeeze}
设 \(A\) 为原始非负矩阵，\(\rho=\rho(A)>1\)，维数 \(d:=\dim A\)，并以命题 \ref{prop:finite-part-residue-reduced-determinant} 的记号定义主极点留数常数 \(C(A)\)。
若核版 RH（平方根界）成立，即对一切非 Perron 特征值 \(\mu\neq\rho\) 有 \(|\mu|\le \rho^{1/2}\)，则有双侧界
\[
(1+\rho^{-1/2})^{-(d-1)}\ \le\ C(A)\ \le\ (1-\rho^{-1/2})^{-(d-1)}.
\]
特别地，当 \(\rho\ge 4\) 时（从而 \(\rho^{-1/2}\le \tfrac12\)），有粗界
\[
|\log C(A)|\ \le\ 2(d-1)\rho^{-1/2}.
\]
\end{proposition}

\begin{proof}
由命题 \ref{prop:finite-part-residue-reduced-determinant}，
\(
C(A)=\prod_{j=2}^{d}(1-\mu_j/\rho)^{-1}
\)（\(\mu_1=\rho\)）。
原始非负矩阵的谱关于共轭闭合，故 \(C(A)>0\) 且
\(
C(A)=\prod_{j=2}^{d}|1-\mu_j/\rho|^{-1}
\)。
在 \(|\mu_j|/\rho\le \rho^{-1/2}\) 下有
\(
1-\rho^{-1/2}\le |1-\mu_j/\rho|\le 1+\rho^{-1/2}
\)，逐项相乘即得双侧界。
当 \(\rho^{-1/2}\le\tfrac12\) 时，用 \(\log(1-x)^{-1}\le 2x\)（\(0<x\le\tfrac12\)）与 \(\log(1+x)\le x\le 2x\) 得
\[
-(d-1)\log(1+\rho^{-1/2})
\ \le\ \log C(A)\ \le\ (d-1)\log(1-\rho^{-1/2})^{-1},
\]
从而 \(|\log C(A)|\le 2(d-1)\rho^{-1/2}\)。证毕。
\end{proof}

\begin{theorem}[异常—RH—维数三岔：线性留数漂移迫使 RH 破裂或维数指数爆炸]\label{thm:finite-part-residue-drift-trichotomy}
对某个 moment-kernel 族 \(\{A^{(q)}_\theta\}_{q\ge 2}\)（\(\theta\) 固定），记
\(\rho_q:=\rho(A^{(q)}_\theta)\)、\(d_q:=\dim A^{(q)}_\theta\) 与 \(C_q(\theta):=C(A^{(q)}_\theta)\)。
假设存在常数 \(c>1\) 使 \(\rho_q\ge c^q\) 对所有充分大 \(q\) 成立，并假设维数次指数增长
\[
\log d_q=o(q)\qquad(q\to\infty).
\]
若核版 RH 对所有充分大 \(q\) 成立，则必有
\[
\log C_q(\theta)=o(q)\qquad(q\to\infty).
\]
因此，任一满足
\[
\liminf_{q\to\infty}\frac{1}{q}\,|\log C_q(\theta)|>0
\]
的“线性留数漂移”现象，都必然迫使下述二者之一发生：
\begin{enumerate}
  \item 存在某个有限 \(q\) 使核版 RH 失效；
  \item \(d_q\) 不能保持次指数增长，必须进入指数增长区间。
\end{enumerate}
\end{theorem}

\begin{proof}
由命题 \ref{prop:finite-part-residue-constant-rh-squeeze}，当 \(q\) 充分大使 \(\rho_q\ge 4\) 时有
\(
|\log C_q(\theta)|\le 2(d_q-1)\rho_q^{-1/2}
\le 2d_q\,c^{-q/2}
\)。
在 \(\log d_q=o(q)\) 下，右端为 \(o(q)\)，故 \(\log C_q(\theta)=o(q)\)。
若 \(\liminf_{q\to\infty}\frac{1}{q}|\log C_q(\theta)|>0\) 仍成立，则与上述 \(o(q)\) 结论矛盾；因此只能落入 (1) 或 (2)。证毕。
\end{proof}

\begin{corollary}[预言机—异常等价律：非零异常必以“有限见证”或“维数爆炸”出现]\label{cor:finite-part-oracle-anomaly-equivalence}
在本文 \(\mathbf{PW}\) 语义中，设某个“异常量”在 \(q\)-重 moment 放大下呈线性规模，并且其谱侧见证可归约为留数常数漂移的线性下界，即存在常数 \(\alpha>0\) 使
\[
|\mathsf{Anom}(q)|\not\equiv 0
\quad\Longrightarrow\quad
\liminf_{q\to\infty}\frac{1}{q}\,|\log C_q(\theta)|\ge \alpha.
\]
则若该异常量非零，必有且仅有如下二者之一发生：存在有限阶 \(q\) 的核版 RH 反例见证；或 \(d_q\) 至少指数增长，从而无限核性质只能以预言机式的无限族检查（或外置模式轴）表达。
\end{corollary}

\begin{proof}
由假设，异常非零迫使 \(|\log C_q(\theta)|\) 出现线性漂移。
在 \(\rho_q\ge c^q\) 与 \(d_q\) 次指数增长的常规尺度假设下，定理 \ref{thm:finite-part-residue-drift-trichotomy} 立即给出二择一结局。证毕。
\end{proof}

